% Reconstructed after the 2026-11-19 10:14 UTC scratch reset.
% Auxiliary ingredients of the independently audited C.3aq; no new result ID.
\paragraph{A quotient of relative $K_2$ in every characteristic.}
Let $S=R\oplus I$ be a split square-zero extension of commutative rings.
There is a natural surjection
\[
                \tau:K_2(S,I)\longrightarrow I\otimes_R\Omega^1_{R/\mathbb Z}
\]
specified, in Weibel's modern Dennis--Stein convention, by
\[
 \tau\langle m,r\rangle=-m\,d\bar r,\qquad
 \tau\langle r,m\rangle=m\,d\bar r\qquad(m\in I).
\]
In particular symbols with both entries in $I$ have zero image.
The relative presentation is
\citet[Chapter III, Theorem 5.11.1(b), p.~44]{frWeibelKbookIII}.
Its hypothesis that $I$ is a radical ideal means that $I$ is contained
in the Jacobson radical, as in that chapter; a square-zero ideal qualifies.
The relation (D1) gives skew symmetry. In (D2), the correction term
$rst$ disappears after multiplication by the entry in $I$, or after
reduction to $R$, because $I^2=0$. Relation (D3) is the Leibniz rule.
For example, if its middle entry is $m\in I$, its two sides map to
$mt\,dr$ and $-rm\,dt+m\,d(tr)$, respectively. The cases with the
first or third entry in $I$ give the same rule; multiple entries in $I$
give zero. Thus all relations hold, and the elements $m\,dr$ show
surjectivity. This is the negative of the map in Weibel's
Exercise 5.14(a). It gives the positive Steinberg normalization
\[
                          \tau\{1+tc,u\}=tc\,d\log u.
\]
In characteristic two we claim a quotient, not an isomorphism.
No surjectivity on sheaf cohomology follows merely from this local
surjection. Naturality and homogeneity in $I$ make the map glue when
$I=L$ is a line bundle. For a perfect field of positive characteristic,
the resulting target is $L\otimes\Omega^1_{V/k}$.

\paragraph{Uniform smooth surface data for every prime.}
Fix a prime $p$, put
\[
 e=p+1,\qquad n=p(p+1),\qquad d=n-3>0,
\]
and work over an algebraically closed field of characteristic $p$.
Consider the smooth plane curve
\[
 C:\quad Y^n-YW^{n-1}=Z^{n-1}W.
\]
It has the unique point $\infty=(0:1:0)$ at infinity, and for
$y=Y/W$, $z=Z/W$ we have
$\operatorname{div}(dz)=n(n-3)\infty$.
Set $D_1=ed\,\infty$, so that $\operatorname{div}(dz)=pD_1$.
To verify the exact local normalization, put $q=Y/Z$ and $h=W/Z$.
Then
\[
 h=q^n-qh^{n-1},\qquad
 z-q^{-n}=q^{n(n-3)}c,\qquad
 c=q(h/q^n)^{n-2}.
\]
The quotient $h/q^n$ is a unit with residue one, and $c$ is a
uniformizer. In particular $dc$ is a unit differential.

The Frobenius-killed extension constructed from $dz$ is
$0\to\mathcal O_C(-D_1)\to E\to\mathcal O_C\to0$.
Use the lines convention for $\mathbb P(E)$; its section $S_\infty$
has normal bundle $\mathcal O_C(D_1)$. The inseparable curve $G$ is
locally $s^p=z$, is smooth and disjoint from $S_\infty$, and satisfies
\[
                    G\sim pS_\infty-pg^*D_1.
\]
This minus sign is obtained from the local equations and the normal
bundle, rather than from the conflicting display in RIMS1736, p.~5.
The cyclic cover of degree $e=p+1$ has inverse-root line bundle
\[
                   M=\mathcal O(-S_\infty+pd\,g^*\infty)
\]
and branch divisor $G+S_\infty$. Affinely it is
$T^{p+1}=s^p-z$. The branch components are smooth and disjoint,
and $e$ is prime to $p$, so the cover is smooth. Valuation one at
a branch component proves that it is connected even when $e$ is composite.
Let $S$ be the reduced section above $S_\infty$.
Then $N_{S/V}=\mathcal O_C(d\infty)$ and $S^2=d$.
Put
\[
                     D=pS+d\,g^*\infty,
                     \qquad L=pD.
\]
We have $D^2=(e^2-1)d$, $D\cdot S=ed$ and $D\cdot F=p$.
Each fibre is integral (its cusp has coprime exponents $p$ and $p+1$).
Every horizontal curve other than $S$ has nonnegative intersection with
$S$ and positive intersection with a fibre. Thus Nakai--Moishezon
proves that $D$, and therefore $L$, is ample.

The local change at base infinity is
$s=q^{ed}s_i+q^{-e}$ and $T=q^{pd}T_i$.
At $S$, set $x=1/s_i$ and $w=T_i/s_i$. The cover equation becomes
\[
                         w^{p+1}=x-cx^{p+1},
 \qquad dT_i=-\frac{x^p}{w^p}\,dc.
\]
Consequently the order of $dT_i$ at $S$ is $p^2$; the base-infinity
factor contributes $pd$. The local unit calculation also checks the
intersection point of these divisors. We obtain
\[
                    \operatorname{div}(dT)=pD,
 \qquad 0\ne\eta=dT\in H^0(V,\Omega^1_V\otimes L^{-1}).
\]
These are the uniform surface ingredients used in C.3aq. They specialize
to valid data in characteristics two and three; C.3an uses a smaller
separate characteristic-five example.

\paragraph{Actual curve germs, without an irreducibility assumption.}
In a trivialization of $L$ with fibre coordinate $t$, the equations
\[
 s=0,\quad T=at,\quad z=-a^{p+1}t^{p+1},\quad
 y^n-y=z^{n-1}
\]
define a curve smooth at $(y,t)=(0,0)$ because the derivative with
respect to $y$ is $-1$. We take the unique irreducible component through
that smooth point. We make no claim that the entire displayed
polynomial is irreducible. Its closure meets the exceptional surface
in finitely many points; normalization over its cone image is finite.
After strict henselization, the chosen branch is a finite local factor
and is already normal at its selected closed point. This produces an
actual integral curve on the fixed strict-local resolution. The symbol
calculation of the preceding lemma gives residue $a$, up to the fixed
simultaneous orientation convention. The global Chow construction and
the quotient by ordinary algebraic equivalence are proved separately
in C.3aq.
